\scanpage{29}
\mspage{13}
% archive sheet number: 28
correspondantes de relèvements \inserted{de $\{1,\dot\tau\}$} sont resp. les relèvements en $\tau_0$, en $\tau_1$ et en $\tau_\infty$ --- qui ne donnent rien de nouveau. Reste à considérer $\tau\dot\sigma_1$ (et ses conjugués \inserted{par $\rho$} $\tau\dot\sigma_\infty$, $\unsure{\tau\dot\sigma_0|\tau_\infty\dot\sigma_0}$), on a $\dot\sigma_1(2)=1/2$, donc
\[
(\tau\dot\sigma_1)(2)=\frac12
\]
et
\[
U^{\tau\dot\sigma_1}=\{z\in U\mid z\overline z=1\}=\mathbb U-\{1\},
\]
où $\mathbb U=\{z\in\mathbb C\mid |z|=1\}$. Cet ensemble contient $P$ et $\overline P$, et ne donne donc pas de \unsure{nouveaux|nouveau} \unsure{relèvements|relèvement} \unsure{p.v.|p.r.} à celui déjà donné, avec $\widetilde\sigma_1$ relevant $\dot\sigma_1\tau$.

Il faudrait quand-même donner aussi une description de $\unsure{\mathscr E|\mathcal E}$ \unsure{tt|tout} entier par gén. et relations, en partant d'un syst. de gén. de $\Gamma=\unsure{\mathfrak S_3|\mathscr S_3|G_3}\times\mathbb Z/2$, qu'on va relever de façon judicieuse. \deleted{On prendra les \inserted{quatre} générateurs}
\[
\deleted{\unsure{\rho'|\rho}=\rho\dot\tau,\quad
\underset{\dot\sigma_0\tau}{\overset{\widetilde{\dot\sigma}_0}{\text{''}}},\quad
\underset{\dot\sigma_1\tau}{\overset{\widetilde{\dot\sigma}_1}{\text{''}}},\quad
\underset{\dot\sigma_\infty\tau}{\overset{\widetilde{\dot\sigma}_\infty}{\text{''}}}}
\]
\deleted{\inserted{(donc $\rho'^{\,6}=1$, $\rho'^{\,3}=\dot\tau$, $\rho'^{\,2}=\rho^2=\rho^{-1}$)}}
\deleted{avec les relations}
\[
\deleted{\left\{\begin{aligned}
\widetilde{\dot\sigma}_1^2&=\widetilde{\dot\sigma}_2^2=\widetilde{\dot\sigma}_3^2=1\\
\widetilde{\dot\sigma}_0\widetilde{\dot\sigma}_1&=\widetilde{\dot\sigma}_1\widetilde{\dot\sigma}_\infty=\widetilde{\dot\sigma}_\infty\widetilde{\dot\sigma}_0=\unsure{\rho'^{\,2}|\rho^2}\\
\rho'^{\,6}&=1\\
\rho'^{\,3}&\text{ commute aux }\dot\sigma_i
\end{aligned}\right.}
\]
